\documentclass[11pt]{article}
\usepackage[margin=0.9in]{geometry}
\usepackage{amsmath,amssymb,amsthm,booktabs,mathrsfs}
\usepackage[T1]{fontenc}
\usepackage{lmodern,microtype}
\usepackage[colorlinks=true,linkcolor=blue,urlcolor=blue]{hyperref}
\newtheorem{theorem}{Theorem}[section]
\newtheorem{lemma}[theorem]{Lemma}
\newtheorem{proposition}[theorem]{Proposition}
\newtheorem{corollary}[theorem]{Corollary}
\theoremstyle{remark}
\newtheorem{remark}[theorem]{Remark}
\newcommand{\norm}[1]{\left\|#1\right\|}
\title{Toward a four-dimensional Yang--Mills construction\\
\large Notebook 22: transport into the next coarse background}
\author{Research working document}
\date{17 October 2026}
\setlength{\emergencystretch}{5em}
\begin{document}
\maketitle
\begin{abstract}
We construct a strictly supported expansion of the actual physical
stationary current background as a function of the next-coarse
field, uniformly in preceding classical refinement and periodic
volume. An affine version of the exact ordinary pivot chart and
the retained flat inverse yield a support-valued stationary
equation. Coordinatewise Cauchy estimates convert the previously
localized scalar action into an analytic gradient majorant without
a factorial loss. The solution is identified with the established
gauge-zero classical branch. We then transport holomorphic
connected-carrier activities through that branch with explicit
losses of analytic radius and support weight. The number of
Taylor coordinates is counted exactly, and the vacuum term is
separated. This supplies a bounded one-step transport for the
generated quantum activities. It does not assert contraction
or stability of the analytic domains under repeated integration.
\end{abstract}

Notebook 21 localized the classical action and controlled one
integration of it in independent current backgrounds. A new
coarse effective action requires the substitution of its
stationary current background. The following two results prove
that substitution and its activity norm estimate. Their domain
and weight losses are retained as inputs to the remaining
quantum induction.


\section{Strict support of the physical stationary background}
\label{st-section}

Fix odd $L$, dimension four, color and the compatible periodic source
geometry with distinct block interfaces. Current spacing is one, old
fine spacing is $M^{-1}$ with $M=L^k$, and next-coarse spacing is $L$.
Write $V=e^{ia}$ and $W=e^{iw}$. These $a,w$ are link angles. In current
physical units the next-coarse potential is $w/L$; in next-coarse unit
coordinates it is $w$. All formulas below use $w$, not $w/L$.
Constants may depend on fixed $L$, but not on $M$ or volume.

Group current coordinates into $L$-block cells, indexed by next-coarse
cells. The retained nonpivot coordinate $z$ and next-coarse datum $w$
have fixed real dimensions per block
\begin{equation}
 D_z=d(L^d-1)\dim\mathfrak g,\qquad D_w=d\dim\mathfrak g,
 \qquad D=D_z+D_w=dL^d\dim\mathfrak g.                         \label{st-dimensions}
\end{equation}
No fluctuation coupling has yet been inserted in $z$. Fixed-dimensional
norm equivalences are included in all radii. Scalar Cauchy estimates use
an actual coordinate product polydisc after this conversion.

We use the previously established ordinary source average and pivot
chart, the physical gauge-zero branch obtained by the two minimum
comparisons, and the strict scalar classical-action expansion. The flat
retained precision has the uniform bounds
\begin{equation}
 A_{M,0}=K_L^T(q_M+T_L^TT_L)K_L,\qquad C_{M,0}=A_{M,0}^{-1},
 \qquad \gamma_*I\le A_{M,0}\le B_*I,                         \label{st-flat-precision}
\end{equation}
and both matrices have uniformly exponentially summable block rows and
columns. Strict scalar support is available with any prescribed finite
cardinality reserve, at the cost of a smaller uniform product radius.
The result proved here is one uniform classical-background transport,
not a claim of closure of repeated quantum RG steps.

\subsection{A local affine pivot chart}
\label{st-affine-section}

At the identity background let
\begin{equation}
 F_1(a)=\frac1i\log M_L(e^{ia})=A_1a+C_1(a),\qquad
 C_1(0)=DC_1(0)=0.                                           \label{st-constraint}
\end{equation}
The actual source normalization is
\begin{equation}
 A_1=LQ_L^{\rm flat},\qquad A_1h_1=I,\qquad A_1K_L=0,
 \qquad h_1(c)=L^{d-1}I\text{ at }p(c).                       \label{st-pivot-normalization}
\end{equation}
The pivot coefficient in $Q_L^{\rm flat}$ is $L^{-d}$, hence its
coefficient in $A_1$ is $L^{1-d}$. Nonpivot restriction of $K_L$ is the
identity; $h_1$ has only pivot entries.

Prescribe the next-coarse datum $w$ and the retained absolute logarithmic
coordinates $z$ by solving
\begin{equation}
 a=K_Lz+h_1w-h_1C_1(a).                                     \label{st-affine-equation}
\end{equation}
Each ordinary constraint equation reads its own pivot and no other
pivot. Thus this is a family of independent finite-dimensional
Lie-algebra-block implicit equations, with neighboring nonpivot values
as parameters. Their derivative in the pivot at zero is invertible.
There are only finitely many local types. A common small contraction
therefore gives a holomorphic solution $a=\mathfrak a(w,z)$ satisfying
\begin{align}
 F_1(\mathfrak a(w,z))&=w,&
 \mathfrak a(w,z)|_{\rm nonpivot}&=z,\notag\\
 \|\mathfrak a(w,z)\|_\infty
    &\le C_L(\|w\|_\infty+\|z\|_\infty),&
 \mathfrak a(w,z)&=K_Lz+h_1w+O((\|w\|+\|z\|)^2).
                                                               \label{st-affine-bounds}
\end{align}
Iteration does not spread to another pivot equation. The dependency
neighborhood and all local analytic constants are fixed at $L$ and
independent of $M$ and volume. This is the unique small constraint
solution with the specified nonpivot values, equivalently the
identity-background ambient source straightening evaluated at
$K_Lz+h_1w$.

\subsection{Scalar support and the exact flat mixed term}
\label{st-scalar-section}

Let $G$ be the actual averaged-contour exponential gauge action, in the
coupling-free normalization, and define
\begin{equation}
 T_M(w,z)=S_M(\mathfrak a(w,z))+G(e^{i\mathfrak a(w,z)}).
                                                               \label{st-action}
\end{equation}
The strict scalar support theorem applies to this fixed-range analytic
substitution and the local gauge action. After fixed $L$-block
regrouping it yields
\begin{align}
 T_M(w,z)&=\sum_{X\ne\varnothing}T_X(w|_X,z|_X),             \label{st-scalar-expansion}\\
 \mathcal K&:=\sup_i\sum_{X\ni i}
 e^{\sigma\tau(X)+A_{\mathrm c}|X|}\|T_X\|_\rho<\infty.         \label{st-scalar-norm}
\end{align}
Here $\tau$ is graph Steiner length on the $L$-block graph, allowing
auxiliary vertices, $\sigma>0$, and $A_{\mathrm c}>0$ is any prescribed
finite reserve. The coordinate-polydisc radius $\rho$ and $\mathcal K$
are uniform in $M$ and volume. Every coefficient starts at degree at
least two: the action does so, the chart vanishes at zero, and the gauge
action has no constant or linear part at the identity.

Put $P_M=q_M+T_L^TT_L$. The full scalar action is stationary at zero, so
only the linear chart enters its quadratic term:
\begin{equation}
 T_M^{(2)}(w,z)=\tfrac12
 \langle K_Lz+h_1w,P_M(K_Lz+h_1w)\rangle.                    \label{st-quadratic-action}
\end{equation}
Thus
\begin{align}
 D_zT_M(w,z)&=A_{M,0}z+B_Mw+N_M(w,z),\notag\\
 B_M&=K_L^TP_Mh_1=K_L^Tq_Mh_1,&
 N_M(0,0)&=DN_M(0,0)=0.                                    \label{st-gradient-split}
\end{align}
The last simplification follows from $T_Lh_1=0$, because the root
contours are block internal and contain no pivot. The rectangular
$B_M$ has uniform exponential block bounds. The pairings here are
counting pairings of the stated angle coordinates; no additional
$L^4$ transpose factor is introduced. Replacing $w$ by the next-coarse
physical potential means substituting $w=L a_{\rm next}$ explicitly.

\subsection{An analytic mixed-slot bound from scalar supports}
\label{st-gradient-section}

The supports in \eqref{st-scalar-expansion} are not uniformly bounded.
The finite-range vertex lemma therefore cannot be applied directly to
their gradient. In particular, crude polarization followed by
$|X|^n$ slot counting can produce a factorial loss. Coordinatewise
Cauchy gives the needed stronger conclusion.

Write $u=(w,z)$ in a fixed scalar basis. On the coordinate polydisc
there are $D|X|$ variables. If $M_X=\|T_X\|_\rho$, its scalar Taylor
coefficients satisfy
\begin{equation}
 |c_{X,\beta}|\le M_X\rho^{-|\beta|}.                       \label{st-scalar-cauchy}
\end{equation}
For one retained derivative coordinate $\alpha$ at a cell $i\in X$,
differentiating the product of geometric series gives, for $t<\rho$,
\begin{equation}
 \sum_{\beta\ge e_\alpha}
 \beta_\alpha |c_{X,\beta}|t^{|\beta|-1}
 \le\frac{M_X}{\rho}(1-t/\rho)^{-D|X|-1}.                   \label{st-gradient-cauchy}
\end{equation}
Sum the at most $D_z$ output components. A symmetric multilinear Taylor
coefficient distributes each scalar monomial coefficient among its
ordered slots with multinomial weights summing to one. Consequently
the sum of its ordered cell-slot operator norms is bounded by the same
absolute scalar coefficient sum, with fixed norm-conversion constants
absorbed in this bound. Every output and input-slot root lies in $X$,
so their rooted Steiner length is at most $\tau(X)$.

For $0<\kappa\le\sigma$, the supremum of the all-order positive row
series is therefore bounded by
\begin{equation}
 \mathcal G(t)=\frac{D_z}{\rho}\sup_i\sum_{X\ni i}
 e^{\kappa\tau(X)}M_X(1-t/\rho)^{-D|X|-1}.                  \label{st-row-series}
\end{equation}
Choose $t_0<\rho/2$ such that
\begin{equation}
 D\log\frac1{1-t_0/\rho}\le\frac{A_{\mathrm c}}2.
 \qquad\text{Then}\qquad
 \mathcal G(t_0)\le\frac{2D_z\mathcal K}{\rho}=:M_g.         \label{st-majorant-radius}
\end{equation}
The supremum over output cells cannot be interchanged with the order
sum. Instead let $v_n$ be the supremum of the order-$n$ mixed-slot row
norm. Equation \eqref{st-row-series} gives the separate bound
$v_n\le M_gt_0^{-n}$. Removing degrees zero and one gives the positive
majorant for $N_M$:
\begin{equation}
 \Phi(t)=\sum_{n\ge2}v_nt^n
 \le\frac{M_g(t/t_0)^2}{1-t/t_0},\qquad
 \Phi'(t)\le\frac{8M_gt}{t_0^2}\quad(0\le t\le t_0/2).
                                                               \label{st-nonlinear-majorant}
\end{equation}
This deliberate geometric factor makes the order/root-supremum
bookkeeping valid. It leaves an analytic quadratic majorant uniform in
refinement and volume.

The original $X$ is a fixed coefficient carrier at the flat baseline,
not an additional variable support. After substituting families for
$w,z$, their union is the actual variable support. Internal rooted
connections through $X$ have already been paid by
\eqref{st-row-series}. This preserves strict dependence on the original
next-coarse variables.

\subsection{The support-valued stationary fixed point}
\label{st-contraction-section}

For a prescribed finite $a_{\mathrm c}\ge0$, use
\begin{equation}
 \|u\|_{\kappa,a_{\mathrm c};r}
 =\sup_i\sum_X e^{\kappa\tau(\{i\}\cup X)+a_{\mathrm c}|X|}
             \sup_{\|w|_X\|<r}\|u_{i;X}(w|_X)\|.            \label{st-rooted-norm}
\end{equation}
The source $w$ has singleton norm $e^{a_{\mathrm c}}r$. For joint insertion
in $N_M$, use the maximum of the separate $w$- and $z$-family norms and
color slots by scalar coordinate and field type. The positive scalar
majorant above already sums all such colors and types, while each
scalar-component family is bounded by its vector-family norm. This
does not identify that maximum with
$\sum_X\|(w_X,z_X)\|$, which could be larger.

Linear exponential block bounds act in \eqref{st-rooted-norm} by their
weighted row sums. The rooted Steiner union inequality and
$e^{a_{\mathrm c}|\cup X_r|}\le\prod_r e^{a_{\mathrm c}|X_r|}$ allow the
proved support-substitution lemma to use
\eqref{st-nonlinear-majorant}. Stationarity becomes
\begin{equation}
 z=-C_{M,0}B_Mw-C_{M,0}N_M(w,z).                            \label{st-stationary-fixed-point}
\end{equation}
Let $c_0$ bound the weighted norm of $C_{M,0}$ and choose
$\ell_0\ge\max\{1,\|C_{M,0}B_M\|_\kappa\}$ uniformly. Put
\begin{equation}
 r_*\le\min\left\{\frac{t_0}{2},
                  \frac{t_0^2}{16c_0M_g}\right\},\qquad
 r\le\frac{r_*}{2\ell_0e^{a_{\mathrm c}}}.                     \label{st-data-radius}
\end{equation}
Decrease these radii further so that reconstruction and evaluation stay
within the common prior chart and uniqueness neighborhoods. On the
family ball $\|z\|\le2\ell_0e^{a_{\mathrm c}}r$, the map
\eqref{st-stationary-fixed-point} is a contraction with constant at most
$1/2$ and preserves that ball. Hence it has a unique holomorphic
supported solution with
\begin{equation}
 \|z_*\|_{\kappa,a_{\mathrm c};r}\le2\ell_0e^{a_{\mathrm c}}r,
 \qquad
 \|z_*+C_{M,0}B_Mw\|_{\kappa,a_{\mathrm c};r}\le C_{a_{\mathrm c}}r^2.
                                                               \label{st-retained-support}
\end{equation}
Evaluation gives a genuine stationary solution on each finite
regulator. Picard iteration defines an absolutely convergent rooted
expansion with original $w$-cell leaves.

The finite-range chart $\mathfrak a$ can now be substituted by the
finite-range support lemma. Define
\begin{equation}
 a_*(w)=\mathfrak a(w,z_*(w)),\qquad
 H_{{\rm cur},M}=h_1-K_LC_{M,0}B_M.                          \label{st-current-extension}
\end{equation}
Then
\begin{equation}
 \|a_*\|_{\kappa,a_{\mathrm c};r}\le C_{a_{\mathrm c}}r,
 \qquad
 \|a_*-H_{{\rm cur},M}w\|_{\kappa,a_{\mathrm c};r}
                                      \le C_{a_{\mathrm c}}r^2.          \label{st-current-support}
\end{equation}
This $H_{{\rm cur},M}$ is the current-angle extension, distinct from the
earlier original-fine extension. It has a uniform exponential
rectangular block bound and $A_1H_{{\rm cur},M}=I$. Its source gauge
condition will follow by identifying the physical branch below; no
gauge-zero conclusion is inferred from stationarity alone. Local
exponentiation also supplies strict families for $e^{ia_*}-I$ and its
inverse minus the identity. For activity transport the bound
$s(r)=C_{a_{\mathrm c}}r$ on the potential itself is sufficient.

\subsection{Identification of the physical gauge-zero branch}
\label{st-identification-section}

For real small $w$, the already proved two-scale minimum comparison and
rooted averaged-contour gauge construction give a branch
$V_{{\rm old},*}(e^{iw})=e^{ia_{{\rm old},*}}$ with
\begin{align}
 M_L(V_{{\rm old},*})&=e^{iw},& P_{V_{{\rm old},*}}(y,x)&=1,\notag\\
 \|a_{{\rm old},*}\|_\infty&\le C_L\|w\|_\infty,&
 S_M(a)&\ge S_M(a_{{\rm old},*})                           \label{st-known-minimum}
\end{align}
for every real feasible $a$ in the common small comparison ball. The
actual gauge action is nonnegative on real unitary fields and vanishes
at $P=1$. Hence
\begin{equation}
 S_M(a)+G(e^{ia})\ge S_M(a_{{\rm old},*})+G(V_{{\rm old},*})
                                                               \label{st-action-comparison}
\end{equation}
on that feasible neighborhood. The affine chart parametrizes every
sufficiently small feasible field by its retained absolute coordinates.
Thus $z_{{\rm old},*}=a_{{\rm old},*}|_{\rm nonpivot}$ is a critical
point of $T_M(w,\cdot)$.

The scalar support expansion and Cauchy with fixed margins give
\begin{equation}
 \|D_z^2T_M(w,z)-A_{M,0}\|_\kappa
                \le C_L(\|w\|_\infty+\|z\|_\infty).         \label{st-convexity}
\end{equation}
For each coefficient, contract its third derivative with the bounded
direction $(w,z)$; summing the second retained matrix index costs only
$|X|$, paid by a cardinality spare. Its input/output distance is at
most $\tau(X)$. Symmetry gives the column bound. This calculation
differentiates the full scalar action, so both chart curvature and
actual gauge curvature remain included. It does not use the ordinary
rooted average's physical adjoint.

Choose a common real retained supremum ball and a data ball making
\eqref{st-convexity} at most $\gamma_*/2$. The retained Hessian is
positive on this convex ball, so there is at most one critical point.
Both $z_*$ from \eqref{st-stationary-fixed-point} and
$z_{{\rm old},*}$ lie inside after a uniform shrink. Therefore
\begin{equation}
 e^{ia_*(w)}=V_{{\rm old},*}(e^{iw}),\qquad
 M_L(e^{ia_*(w)})=e^{iw},\qquad P_{e^{ia_*(w)}}(y,x)=1.       \label{st-physical-identities}
\end{equation}
This identifies with a known gauge-zero minimum; a generic critical
equation was not used to assert the gauge condition. Holomorphy and the
finite-dimensional identity theorem extend these identities from the
real ball to the common complex product domain. Real minimum
inequalities are not extended as inequalities. Linearizing the last
identity gives $T_LH_{{\rm cur},M}=0$.

\subsection{Transport conclusion and its domain limitation}
\label{st-scope-section}

For any prescribed finite cardinality reserve and sufficiently small
positive $\kappa$, \eqref{st-current-support} is now a strict-support
expansion of the actual physical stationary background in its original
next-coarse angle variables. The radius is uniform in old refinement
and volume. Output roots are $L$-blocks of current fields and input
labels are next-coarse cells on the same block graph. If an activity
uses another current-cell convention, its projection to this graph and
the finite fiber multiplicity must be included explicitly; fixed
regrouping is not a contraction estimate. Coefficients can depend on
the fixed baseline and regulator. Strict variable support does not by
itself assert volume consistency of a chosen decomposition.

There is a concrete obstruction to inferring an $L^{-1}$ contraction
of the angle supremum domain. Take small coarse pure-gauge data
$W(y,y')=g_yg_{y'}^{-1}$ and set the current-site gauge to $g_y$
throughout each $L$-block. Its current field is the identity internally
and equals $W$ on the corresponding block-boundary links. It has
$P_V=1$, ordinary average $W$ and zero classical action. By the same
minimum comparison and uniqueness it is the physical stationary
representative. Its boundary angle is of order $\|w\|$, not
$\|w\|/L$. The physical conversion $a_{\rm next}=w/L$ does not change
that fact about the angle-domain map.

This theorem supplies one uniform step and an explicit small-domain
bound $s(r)=C_ar$ for substitution in supplied analytic activities.
Preservation of a quantum activity norm, any associated radius loss,
and a compatible common domain under repeated integrations require
separate proofs. No running-coupling, mixed-history, large-field or
interacting continuum conclusion follows here.

% Notebook 22 quantum-activity transport module. No preamble. Prefix qt-.
\section{Strict-support transport of analytic quantum activities}
\label{qt-section}

We now substitute the actual stationary current background into a
supplied family of quantum activities. This is a bounded finite-step
map. Its input has only one fixed positive support-cardinality reserve;
we do not assume an arbitrary reserve for quantum activities. The
number of Taylor coordinates is proportional to the original carrier
size, and we explicitly pay that cost. Constants are uniform in the
old refinement $M$ and compatible periodic volume, with fixed $L$.

\subsection{Carrier geometry and the precise analytic inputs}
\label{qt-input-section}

Let $\Lambda$ be the input cell graph and $\Lambda'$ the next field's
cell graph. A projection $\pi:\Lambda\to\Lambda'$ has fibers of size
at most $F$ and sends every nonempty connected input carrier to a
connected set. For ordinary $L$-cube blocking, $F=L^d$. Identity
projection is allowed. Write $d'$ for graph distance on $\Lambda'$ and
define
\begin{equation}\label{qt-distance-sum}
 S_\delta=\sup_{z\in\Lambda'}\sum_{j\in\Lambda'}e^{-\delta d'(z,j)}
 \le\sum_{j\in\mathbb Z^d}e^{-\delta|j|_1}<\infty,
 \qquad \delta>0.
\end{equation}
The bound is uniform on compatible tori: choose one minimal-distance
integer representative of every periodic site. Short periodic
directions do not enlarge the sum.

Each input cell carries at most $D$ scalar complex background
coordinates. For a nonempty connected $H\subset\Lambda$, suppose
$E_H(a|_H)$ is bounded and holomorphic on
\begin{equation}\label{qt-input-domain}
 |a_i^\alpha|<b,\qquad i\in H,\quad 1\le\alpha\le D,
\end{equation}
and put $e_H=\norm{E_H}_b$. The only activity hypothesis is
\begin{equation}\label{qt-input-norm}
 \sup_{i\in\Lambda}\sum_{H\ni i}e^{a_0|H|}e_H\le\epsilon,
 \qquad a_0>0\text{ fixed}.
\end{equation}
Here $D$ counts background scalar coordinates, not automatically
retained Gaussian coordinates. A product of Euclidean fiber balls of
radius $b_E$ contains the scalar polydisc of radius $b_E/\sqrt D$;
make this fixed norm conversion before applying the formulas below.

The actual stationary potential on the next-field product domain of
radius $r$ has the strict support expansion
\begin{align}
 a_i^\alpha(w)&=\sum_{\varnothing\ne X\subset\Lambda'}
                   u_{i,\alpha;X}(w|_X),
 &u_{i,\alpha;X}(0)&=0,                                  \label{qt-background-series}\\
 \sup_{i,\alpha}\sum_X
 e^{\kappa\tau(\{\pi(i)\}\cup X)}
       \norm{u_{i,\alpha;X}}_r&\le s(r)\le C_a r,
 &\kappa&>0.                                             \label{qt-background-norm}
\end{align}
Here $\tau$ is graph Steiner length, allowing auxiliary vertices.
The preceding stationary-background theorem supplies this estimate;
its stronger support-cardinality weight is not needed here. A
maximum-coordinate bound on its full current-field block fiber
implies \eqref{qt-background-norm} for each scalar component, without
summing over the fiber.

If only the total $a_i(0)$ vanishes, subtract each summand's constant.
Absolute convergence preserves the total and increases $s$ by at most
a factor two. Incorporate that factor into $C_a$. For a nonzero
reference potential, center both \eqref{qt-input-domain} and
\eqref{qt-background-series} at that reference and use the available
distance to the analytic boundary in place of $b$.

\begin{remark}[The present physical indexing]\label{qt-indexing}
The normalized Gaussian activities supplied by Notebooks 18 and 21
are already indexed by $L$-blocks of the current unit lattice. The
next-coarse $w$ labels are the same block index graph. The present
physical substitution therefore uses
\[
 \pi=\mathrm{id},\qquad F=1,\qquad
 D=D_{\mathrm{bg}}=dL^d\dim\mathfrak g.
\]
The last dimension includes all current background links assigned to
the block; it differs from the retained Gaussian dimension
$d(L^d-1)\dim\mathfrak g$. After substitution, these labels are unit
cells of the next current lattice. A further projection for the
following integration is a distinct operation, discussed below.
It must not be counted as an additional gain during this substitution.
\end{remark}

\subsection{Bounded transport with an explicit loss budget}

Choose $a_{\mathrm{out}},\mu\ge0$ and $\delta>0$, and set
\begin{equation}\label{qt-budget-notation}
 \beta=a_{\mathrm{out}}+\mu,\qquad
 t=\frac{s(r)}b<1,\qquad
 \alpha_D(t)=D\log\frac1{1-t}.
\end{equation}
For a connected output carrier $K$, let $\ell(K)$ denote metric
minimum spanning-tree length. Define
\begin{equation}\label{qt-output-norm}
 \norm{E'}_{a_{\mathrm{out}},\mu;r}
 =\sup_{z\in\Lambda'}\sum_{K\ni z}
      e^{a_{\mathrm{out}}|K|+\mu\ell(K)}\norm{E'_K}_r.
\end{equation}

\begin{theorem}[Finite-step activity transport]\label{qt-theorem}
Assume \eqref{qt-input-norm}--\eqref{qt-background-norm} and
\begin{equation}\label{qt-budget}
 \beta+\delta\le\kappa,\qquad
 \beta+\alpha_D(t)\le a_0.
\end{equation}
There are strictly supported holomorphic activities $E'_K(w|_K)$
on nonempty connected $K\subset\Lambda'$, with $E'_K(0)=0$, such that
\begin{align}
 \sum_H E_H(a(w)|_H)
   &=\mathrm{Vac}_\Lambda+\sum_K E'_K(w|_K),
 &\mathrm{Vac}_\Lambda&=\sum_H E_H(0),                 \label{qt-exact-identity}\\
 \norm{E'}_{a_{\mathrm{out}},\mu;r}
   &\le F S_\delta\epsilon.                            \label{qt-transport-bound}
\end{align}
All regroupings converge absolutely. The bound is independent of
$M$ and finite periodic volume. It is not a contraction assertion.
\end{theorem}

A concrete choice with strictly positive output weights is
\begin{align}
 \beta&=\min\{a_0/4,\kappa/4\},
 &a_{\mathrm{out}}&=\mu=\beta/2,
 &\delta&=\kappa/4,\notag\\
 \frac{s(r)}b&\le\min\{1/2,a_0/(8D)\}.                 \label{qt-concrete-budget}
\end{align}
Indeed $\alpha_D(t)\le2Dt\le a_0/4$. Thus a sufficient radius is
\begin{equation}\label{qt-radius}
 r\le\frac b{C_a}\min\{1/2,a_0/(8D)\},
\end{equation}
together with the stationary-background theorem's own radius
restriction. All centering and norm-conversion constants are included
in $b$ or $C_a$. Nothing requires $C_a<1$.

\begin{proof}
Expand each activity in its actual $D|H|$ scalar coordinates:
\begin{equation}\label{qt-taylor}
 E_H(a|_H)=\sum_{\nu\in\mathbb N^{D|H|}}c_{H,\nu}a^\nu,
 \qquad |c_{H,\nu}|\le e_H b^{-|\nu|}.
\end{equation}
Coordinatewise Cauchy on smaller polydiscs, followed by a limit,
gives the coefficient bound. There is no dimension-free absolute
Taylor norm being assumed.

For every scalar occurrence in a monomial $a^\nu$, insert
\eqref{qt-background-series}. Its chosen support $X$ inherits the
root $\pi(i)$ of that scalar coordinate. For each such rooted set,
fix a deterministic minimizing connected tree $T(i,X)$ containing
$\{\pi(i)\}\cup X$. Ties can be resolved by any fixed finite-volume
rule. For a decorated term define
\begin{equation}\label{qt-hull}
 K=\pi(H)\cup\bigcup_{\text{occurrences}}\mathrm{Vert}\,T(i,X),
 \qquad T_{\mathrm{tot}}
       =\sum_{\text{occurrences}}\tau(\{\pi(i)\}\cup X).
\end{equation}
The original projected carrier is retained even if some of its
variables are absent from the monomial. It is connected, and every
attached tree has a root in it. Consequently
\begin{equation}\label{qt-hull-geometry}
 |K|\le|\pi(H)|+T_{\mathrm{tot}},\qquad
 z\in K\ \Longrightarrow\ d'(z,\pi(H))\le T_{\mathrm{tot}}.
\end{equation}
A tree with $q$ edges adds at most $q$ vertices after attachment at
its root. A vertex on an attached tree is at distance at most that
tree's length from $\pi(H)$. Overlaps and periodic wrapping can
only improve these inequalities.

The decorated function reads only $w|_K$. Its true variable set may
be smaller: $K$ is a declared connected carrier, not a claim of
minimal dependence. A nonconstant monomial is a product of factors
vanishing at zero and hence also vanishes at zero. Define $E'_K$ as
the sum of these nonconstant terms having carrier $K$. Set aside
the constant Taylor terms as $\mathrm{Vac}_\Lambda$.

For every connected lattice vertex set $K$,
\begin{equation}\label{qt-connected-mst}
 \ell(K)=|K|-1,
 \qquad e^{a_{\mathrm{out}}|K|+\mu\ell(K)}\le e^{\beta|K|}.
\end{equation}
Indeed the induced connected graph has a unit-edge spanning tree,
while every metric tree has $|K|-1$ edges of distance at least one.
Graph Steiner length is used for the attached rooted sets because
deleting intermediate roots from a terminal metric spanning tree
need not decrease its length. At the output all connecting vertices
are retained, so \eqref{qt-connected-mst} applies exactly.

Put $h=|\pi(H)|$ and $k=|H|$. For a fixed output anchor $z$,
\eqref{qt-hull-geometry} gives
\begin{align}
 \mathbf1_{\{z\in K\}}
 e^{a_{\mathrm{out}}|K|+\mu\ell(K)}
 &\le e^{\beta h}e^{-\delta d'(z,\pi(H))}
       \prod_{\text{occurrences}}
       e^{(\beta+\delta)\tau(\{\pi(i)\}\cup X)}.       \label{qt-anchor-hull}
\end{align}
If $z\notin K$ this is automatic; otherwise insert
$1\le e^{\delta(T_{\mathrm{tot}}-d'(z,\pi(H)))}$.
The first condition in \eqref{qt-budget} lets us sum each scalar
occurrence's support using \eqref{qt-background-norm}, at cost $s(r)$.
The exact sum over scalar multiindices is
\begin{equation}\label{qt-coordinate-count}
 \sum_{\nu\in\mathbb N^{Dk}}t^{|\nu|}
    =(1-t)^{-Dk}=e^{\alpha_D(t)k}.
\end{equation}
Omitting $\nu=0$ replaces the last expression by
$e^{\alpha_D(t)k}-1$. Thus the full contribution from a fixed $H$
is bounded by
\begin{equation}\label{qt-fixed-carrier-bound}
 e_H e^{\beta h+\alpha_D(t)k}
                  e^{-\delta d'(z,\pi(H))}.
\end{equation}
Since $h\le k$, the second condition in \eqref{qt-budget} pays
the two exponential costs from the given reserve $a_0 k$.
Finally,
\begin{equation}\label{qt-reroot}
 e^{-\delta d'(z,\pi(H))}
    \le\sum_{i\in H}e^{-\delta d'(z,\pi(i))}.
\end{equation}
Summing \eqref{qt-fixed-carrier-bound} over $H$ therefore gives
\begin{align}
 \sum_i e^{-\delta d'(z,\pi(i))}
             \sum_{H\ni i}e^{a_0|H|}e_H
 &\le F S_\delta\epsilon,                              \label{qt-final-majorant}
\end{align}
which proves \eqref{qt-transport-bound}.

This positive majorant also justifies every substitution and
regrouping. At finite volume, summing over output anchors bounds
the total absolute series. Normal convergence on smaller product
polydiscs proves holomorphy. No order sum has been interchanged
with an output supremum: for each fixed $H$ the scalar series is
summed exactly in \eqref{qt-coordinate-count}, before taking the
anchored carrier sum. In particular, no crude $|H|^n$ slot count
with a missing factorial is hidden in the argument.
\end{proof}

\subsection{Vacuum and the smaller-radius improvement}

The vacuum has the explicit density assignment
\begin{equation}\label{qt-vacuum-density}
 v_i=\sum_{H\ni i}\frac{E_H(0)}{|H|},\qquad
 \mathrm{Vac}_\Lambda=\sum_i v_i,
 \qquad |v_i|\le\epsilon,
 \qquad |\mathrm{Vac}_\Lambda|\le|\Lambda|\epsilon.
\end{equation}
Grouping these densities into a block costs at most $F$.

\begin{corollary}[Vacuum-subtracted radius bound]\label{qt-small-radius}
If the hypotheses of Theorem~\ref{qt-theorem} hold with
\begin{equation}\label{qt-spare-reserve}
 \rho=a_0-\beta-\alpha_D(t)>0,
\end{equation}
then also
\begin{equation}\label{qt-small-radius-bound}
 \norm{E'}_{a_{\mathrm{out}},\mu;r}
 \le F S_\delta\epsilon\,\alpha_D(t) C_\rho,
 \qquad C_\rho=\sup_{k\ge1}k e^{-\rho k}\le\frac1{e\rho}.
\end{equation}
\end{corollary}
\begin{proof}
For nonconstant terms use
$e^{\alpha_D(t)k}-1\le\alpha_D(t)k e^{\alpha_D(t)k}$
in place of \eqref{qt-coordinate-count}. The spare
$e^{-\rho k}$ pays the factor $k$. The remaining rerooting sum
is unchanged.
\end{proof}
One may take the smaller of \eqref{qt-transport-bound} and
\eqref{qt-small-radius-bound}. This is improved smallness on a
strictly smaller field domain. It does not prove a contraction
at a common radius under repeated RG steps.

\subsection{An optional connected-carrier compression estimate}
\label{qt-compression-section}

\begin{lemma}[Connected carriers under cube projection]\label{qt-compression}
For odd $L\ge5$, put
\begin{equation}\label{qt-compression-constants}
 r_L=(L-1)/2,\qquad s_L=\lfloor(r_L-1)/2\rfloor,
 \qquad k_L=s_L+1.
\end{equation}
If $H$ is connected and $\pi_L$ is $L$-cube projection, then
\begin{equation}\label{qt-compression-bound}
 |\pi_L(H)|\le2^d\max\{1,|H|/k_L\}
           \le2^d+(2^d/k_L)|H|.
\end{equation}
The estimate holds on compatible rectangular tori.
\end{lemma}
\begin{proof}
Choose a maximal subset of $H$ separated by at least $r_L$ in
the periodic supremum metric. The supremum balls of radius
$r_L-1$ around its $J$ centers cover $H$. Each such ball has
side length $2r_L-1<L$ and meets at most $2^d$ coarse blocks,
so $|\pi_L(H)|\le2^d J$. If $J=1$ the result follows.
If $J\ge2$, the graph-distance balls of radius $s_L$ about
the centers are disjoint, since $2s_L<r_L$. For each center,
a path in $H$ to another center contains at least $s_L+1$
distinct vertices before leaving its ball. Hence
$Jk_L\le|H|$. Identifications in short periodic directions
only reduce the block counts; the same separation and path
argument uses periodic distances.
\end{proof}

\begin{corollary}[Transport with a proved carrier gain]\label{qt-compressed-transport}
Suppose the projection in Theorem~\ref{qt-theorem} satisfies
$|\pi(H)|\le c_0+c_1|H|$ on connected carriers. The second
condition in \eqref{qt-budget} may be replaced by
\begin{equation}\label{qt-compressed-budget}
 c_1\beta+\alpha_D(t)\le a_0,
\end{equation}
and the resulting bound is
\begin{equation}\label{qt-compressed-bound}
 \norm{E'}_{a_{\mathrm{out}},\mu;r}
       \le e^{\beta c_0}F S_\delta\epsilon.
\end{equation}
The smaller-radius estimate also holds with the factor
$e^{\beta c_0}$ and
$\rho=a_0-c_1\beta-\alpha_D(t)>0$.
\end{corollary}
\begin{proof}
Keep $h=|\pi(H)|$ in \eqref{qt-fixed-carrier-bound}, and
use $\beta h\le\beta c_0+c_1\beta |H|$. All other parts
of the proof are unchanged.
\end{proof}
Lemma~\ref{qt-compression} provides
$c_0=2^d$, $c_1=2^d/k_L$. A sufficiently large fixed $L$
can therefore help pay stronger cardinality weights. But the
attached background trees must still obey
$\beta+\delta\le\kappa$ in the graph used by the substitution.
A larger spatial exponent cannot be assumed without a separate
regrouping proof and its accompanying cardinality and domain costs.
The constant $e^{\beta2^d}L^d S_\delta$ can be large.

In the actual convention of Remark~\ref{qt-indexing}, current
substitution uses identity projection. Afterwards, its connected
carriers are unit-cell carriers of the next current field $w$.
Regrouping them into the following integration's $L$-block Gaussian
cells is a legitimate new projection. It must include that later
pivot-chart enlargement and background-domain budget; it is not
a second projection hidden in \eqref{qt-exact-identity}.

\subsection{What the estimate does and does not establish}

Connectedness of the input carrier is necessary for this
cardinality-only hypothesis. Without it, an activity reading two
sites at arbitrarily large separation can have fixed cardinality
norm, while any connected output carrier containing both has size
at least their distance plus one. A spatial input norm could pay
this cost, but it is not part of \eqref{qt-input-norm}.

There is no automatic contraction: a constant singleton activity
$E_{\{i\}}=c$ at each old cell gives $L^d c$ after summation in
a full coarse block. Even after vacuum subtraction, an identity
background map and a linear singleton activity have no strict
contraction at the same analytic radius. The gain in
\eqref{qt-small-radius-bound} is explicitly a radius loss.

For example, a Gaussian output reserve $a_0\le1/2$ cannot simply
be re-used as an interaction input demanding exponent $4$.
One must verify a compression/enlargement budget, a suitable
stationary-background exponent, and the next domain. The
stationary branch has a one-step radius with an $L$-dependent
constant; no supremum-domain contraction by $L^{-1}$ has been
proved. These facts leave a genuine induction interface.

Every transported coefficient is strictly a function of its
declared $w$ variables. Its fixed baseline may still depend on
$M$ and the finite regulator, and a deterministic choice of
connecting trees need not give coefficientwise volume consistency.
The uniform bounds alone imply no thermodynamic consistency.

Finally, the theorem transports supplied quantum activities
through the matched classical stationary background. It does
not identify the full interacting quantum stationary point with
that classical branch, absorb inherited interactions into a new
minimization, or control running couplings, large-field histories,
or observable insertions. It supplies an exact finite-step bridge
whose coefficient, support, volume, and analytic-radius costs
are all displayed.

\section{The remaining quantum recurrence}
\label{ct22-next}

The actual stationary substitution is now represented by
strictly supported functions of the next-coarse variables.
The connected activities supplied by Notebook 21, and its
separate normalization terms after a common reduction of
their support exponent, satisfy the input form of the
transport theorem. Thus their sum can be expressed as
local next-coarse activities on the explicitly smaller
analytic ball. This is a completed finite-step interface.

For the specific convention used in Notebook 21, the
integration cells are already the next-coarse block cells.
That substitution uses the identity projection. Projection
into another set of $L$-blocks occurs only when preparing
these activities for the following integration. It cannot
be counted twice.

A repeated quantum construction requires more. The generated
activity norm has a fixed small support exponent, whereas the
Gaussian integration theorem asks for a stronger weighted
input. The connected-carrier estimate explains a possible
geometric gain at the later projection, but its constants
and the new fluctuation-chart enlargement must be included.
Notebook 24 makes that inherited-interaction budget explicit;
Notebook 26 meets it using a fixed auxiliary grouping and
recovers a common positive carrier exponent.
The smaller analytic radius in the transport theorem must
also be replaced by a proved stable domain recursion.
The explicit pure-gauge example above rules out obtaining
that recursion merely by declaring that every link angle
contracts by $L^{-1}$ in this rooted representative.
Notebook 25 constructs a smaller pre-root-gauge representative
and identifies the analytic composition germ for an invariant
total functional, with its remaining activity and integral
domain requirements kept explicit.

Finally, the vacuum, marginal gauge-action coefficient and
irrelevant interactions need their appropriate extractions
and scaling estimates. A bounded norm map is insufficient:
vacuum singleton terms already grow by block volume.
Large-field complements, source insertions and the continuum
reconstruction requirements of Notebook 01 remain part of
the full existence argument.

\begin{thebibliography}{9}
\bibitem{prior22}
\emph{Toward a four-dimensional Yang--Mills construction},
research Notebooks 15--21 in this workspace. The physical
stationary branch, exact current pivot chart, retained
inverse, scalar support construction and connected
integration activities are the principal inputs.
\bibitem{balaban10922}
T.~Ba\l aban,
\emph{Renormalization group approach to lattice gauge field
theories I}, Commun. Math. Phys. \textbf{109} (1987), 249--301,
\href{https://scispace.com/pdf/renormalization-group-approach-to-lattice-gauge-field-3ygai0vxt1.pdf}{original article copy}.
\bibitem{balaban11622}
T.~Ba\l aban,
\emph{Renormalization group approach to lattice gauge field
theories. II. Cluster expansions},
Commun. Math. Phys. \textbf{116} (1988), 1--22,
\href{https://scispace.com/pdf/renormalization-group-approach-to-lattice-gauge-field-4penm4k7hc.pdf}{original article copy}.
\bibitem{source22}
E.~Sussman and collaborators,
\href{https://github.com/ethanwsussman/Yang--Mills-LLM-experiment}{Yang--Mills-LLM-experiment},
deposited raw accounts at commit
\texttt{e392ef4501\allowbreak 3a581664ba\allowbreak b76b908671\allowbreak 70b30fe10e}.
\end{thebibliography}
\end{document}
